% --- Archon physics formalization source begin ---
% archon:physics
% archon:covers PhyXMiniProblems/problem_phyx_mini_0690.lean
% archon:source-report reports/phyx_mini/problem_phyx_mini_0690.source.json

\chapter{Physics problem phyx\_mini\_0690}
\label{ch:PhyXMiniProblems_problem_phyx_mini_0690}

\paragraph{Problem source.}
\#\# Physical scenario

Towns A and B in figure are \textbackslash{}( 80.0 \textbackslash{}text\{ km\} \textbackslash{}) apart. A couple arranges to drive from town A and meet a couple driving from town B at the lake, L. The two couples leave simultaneously and drive for \textbackslash{}( 2.50 \textbackslash{}text\{ h\} \textbackslash{}) in the directions shown. Car 1 has a speed of \textbackslash{}( 90.0 \textbackslash{}text\{ km/h\} \textbackslash{}).

\#\# Question

If the cars arrive simultaneously at the lake, what is the speed of car 2?

\#\# Answer choices

- A: A: \textbackslash{}( 67.1 \textbackslash{}text\{ km/h\} \textbackslash{})

- B: B: \textbackslash{}( 54.6 \textbackslash{}text\{ km/h\} \textbackslash{})

- C: C: \textbackslash{}( 68.8 \textbackslash{}text\{ km/h\} \textbackslash{})

- D: D: \textbackslash{}( 32.7 \textbackslash{}text\{ km/h\} \textbackslash{})

\#\# Figure caption (auxiliary; use the image as primary evidence)

The image depicts a map illustrating two paths from points A and B to point L, located near a body of water represented in blue. 

- **Objects and Labels**: 
  - Two cars, labeled 1 (blue) and 2 (red), are shown on the paths.
  - A black path connects point A to point L, with car 1 positioned along it.
  - Another black path connects point B to point L, with car 2 along it.
  - A dashed line connects A to B, marked as 80.0 km.
  - The angle between the path from A to B and the path from A to L is labeled as 40.0°.

- **Areas**:
  - Both A and B are within highlighted regions on the map, with light orange and yellow shading indicating areas of interest or urban areas, while green indicates general land.
  - A light blue river or stream is visible near these regions.

- **Relationships**:
  - The diagram suggests a navigational or route planning context, emphasizing angles and distances between points.
  
This layout likely serves an illustrative purpose, such as explaining navigation or geometry concepts.

\#\# Dataset metadata

Kinematics; Physical Model Grounding Reasoning, Spatial Relation Reasoning

\paragraph{Recorded answer/context.}
C: C: \textbackslash{}( 68.8 \textbackslash{}text\{ km/h\} \textbackslash{})

\paragraph{Figure/image path.}
/root/proposal\_for\_physic/science-mango/phyx\_mini\_run/phyx\_data/test\_image/690.png

\paragraph{Formalization target.}
create a compiling Lean file with sorry bodies at `PhyXMiniProblems/problem\_phyx\_mini\_0690.lean`. The Lean declarations must preserve the physical quantities, dimensions or dimensional roles, figure labels, governing-law hypotheses, and final relation expressed by this problem.
Use Mathlib/Physlib names found through LeanExplore where available. If a physics API is missing, introduce faithful local abstractions rather than scalar placeholder aliases.

\begin{theorem}[Physics formalization target]
\label{thm:physics:phyx_mini_0690:target}
\lean{PhyXMiniProblems.ProblemPhyXMini0690.carTwoSpeed_is_answerChoiceC}
\uses{def:physics:phyx-mini-0690:phyxminiproblems-problemphyxmini0690-speedinkilometresperhour, def:physics:phyx-mini-0690:phyxminiproblems-problemphyxmini0690-angleinradiansfromdegrees, def:physics:phyx-mini-0690:phyxminiproblems-problemphyxmini0690-lakemeetingsetup, def:physics:phyx-mini-0690:phyxminiproblems-problemphyxmini0690-matchesdrivingscenario, def:physics:phyx-mini-0690:phyxminiproblems-problemphyxmini0690-matchessuppliedfigure, def:physics:phyx-mini-0690:phyxminiproblems-problemphyxmini0690-matchesproblemreadouts, def:physics:phyx-mini-0690:phyxminiproblems-problemphyxmini0690-hasphysicallakemeetingparameters, def:physics:phyx-mini-0690:phyxminiproblems-problemphyxmini0690-satisfiesstraightroutekinematics, def:physics:phyx-mini-0690:phyxminiproblems-problemphyxmini0690-isclosestanswerchoice}
The assigned autoformalize agent should translate this physics problem into Lean declarations in the covered file, with theorem and lemma proof bodies written as `by sorry`.
\end{theorem}
\begin{proof}
This is an autoformalization task, not a proof task. Produce statements that can later be proved by the physics prover without weakening the physical model.
\end{proof}

% --- Archon declaration topology begin ---
\section{Declaration topology}
\begin{definition}[Duration Quantity]
  \label{def:physics:phyx-mini-0690:phyxminiproblems-problemphyxmini0690-durationquantity}
  \lean{PhyXMiniProblems.ProblemPhyXMini0690.DurationQuantity}
  A nonnegative, unit-independent physical duration
\end{definition}
\begin{proof}
  This is the declared model definition of Duration Quantity.
\end{proof}
\begin{definition}[Speed Quantity]
  \label{def:physics:phyx-mini-0690:phyxminiproblems-problemphyxmini0690-speedquantity}
  \lean{PhyXMiniProblems.ProblemPhyXMini0690.SpeedQuantity}
  A nonnegative, unit-independent physical speed
\end{definition}
\begin{proof}
  This is the declared model definition of Speed Quantity.
\end{proof}
\begin{definition}[duration In Hours]
  \label{def:physics:phyx-mini-0690:phyxminiproblems-problemphyxmini0690-durationinhours}
  \lean{PhyXMiniProblems.ProblemPhyXMini0690.durationInHours}
  \uses{def:physics:phyx-mini-0690:phyxminiproblems-problemphyxmini0690-durationquantity}
  Read a physical duration in hours
\end{definition}
\begin{proof}
  This is the declared model definition of duration In Hours.
\end{proof}
\begin{definition}[speed In Kilometres Per Hour]
  \label{def:physics:phyx-mini-0690:phyxminiproblems-problemphyxmini0690-speedinkilometresperhour}
  \lean{PhyXMiniProblems.ProblemPhyXMini0690.speedInKilometresPerHour}
  \uses{def:physics:phyx-mini-0690:phyxminiproblems-problemphyxmini0690-speedquantity}
  Read a physical speed in kilometres per hour
\end{definition}
\begin{proof}
  This is the declared model definition of speed In Kilometres Per Hour.
\end{proof}
\begin{definition}[angle In Radians From Degrees]
  \label{def:physics:phyx-mini-0690:phyxminiproblems-problemphyxmini0690-angleinradiansfromdegrees}
  \lean{PhyXMiniProblems.ProblemPhyXMini0690.angleInRadiansFromDegrees}
  Convert a degree readout to the radian-valued angle used by Mathlib
\end{definition}
\begin{proof}
  This is the declared model definition of angle In Radians From Degrees.
\end{proof}
\begin{definition}[Place Label]
  \label{def:physics:phyx-mini-0690:phyxminiproblems-problemphyxmini0690-placelabel}
  \lean{PhyXMiniProblems.ProblemPhyXMini0690.PlaceLabel}
  Place labels printed on the primary map
\end{definition}
\begin{proof}
  This is the declared model definition of Place Label.
\end{proof}
\begin{definition}[Car Label]
  \label{def:physics:phyx-mini-0690:phyxminiproblems-problemphyxmini0690-carlabel}
  \lean{PhyXMiniProblems.ProblemPhyXMini0690.CarLabel}
  The two labeled cars in the primary map
\end{definition}
\begin{proof}
  This is the declared model definition of Car Label.
\end{proof}
\begin{definition}[Kilometre Map Point]
  \label{def:physics:phyx-mini-0690:phyxminiproblems-problemphyxmini0690-kilometremappoint}
  \lean{PhyXMiniProblems.ProblemPhyXMini0690.KilometreMapPoint}
  A point of the Euclidean map. Its two real coordinates and its metric distance are read in kilometres; this is a coordinate readout, not a physical-quantity alias
\end{definition}
\begin{proof}
  This is the declared model definition of Kilometre Map Point.
\end{proof}
\begin{definition}[Lake Meeting Figure]
  \label{def:physics:phyx-mini-0690:phyxminiproblems-problemphyxmini0690-lakemeetingfigure}
  \lean{PhyXMiniProblems.ProblemPhyXMini0690.LakeMeetingFigure}
  \uses{def:physics:phyx-mini-0690:phyxminiproblems-problemphyxmini0690-placelabel, def:physics:phyx-mini-0690:phyxminiproblems-problemphyxmini0690-carlabel}
  Qualitative and numerical information printed on the supplied map
\end{definition}
\begin{proof}
  This is the declared model definition of Lake Meeting Figure.
\end{proof}
\begin{definition}[Lake Meeting Setup]
  \label{def:physics:phyx-mini-0690:phyxminiproblems-problemphyxmini0690-lakemeetingsetup}
  \lean{PhyXMiniProblems.ProblemPhyXMini0690.LakeMeetingSetup}
  \uses{def:physics:phyx-mini-0690:phyxminiproblems-problemphyxmini0690-durationquantity, def:physics:phyx-mini-0690:phyxminiproblems-problemphyxmini0690-speedquantity, def:physics:phyx-mini-0690:phyxminiproblems-problemphyxmini0690-placelabel, def:physics:phyx-mini-0690:phyxminiproblems-problemphyxmini0690-carlabel, def:physics:phyx-mini-0690:phyxminiproblems-problemphyxmini0690-kilometremappoint, def:physics:phyx-mini-0690:phyxminiproblems-problemphyxmini0690-lakemeetingfigure}
  All independent quantities in the problem. In particular, car 2's speed is a free physical observable; it is not defined from an answer choice
\end{definition}
\begin{proof}
  This is the declared model definition of Lake Meeting Setup.
\end{proof}
\begin{definition}[map Distance In Kilometres]
  \label{def:physics:phyx-mini-0690:phyxminiproblems-problemphyxmini0690-mapdistanceinkilometres}
  \lean{PhyXMiniProblems.ProblemPhyXMini0690.mapDistanceInKilometres}
  \uses{def:physics:phyx-mini-0690:phyxminiproblems-problemphyxmini0690-placelabel, def:physics:phyx-mini-0690:phyxminiproblems-problemphyxmini0690-lakemeetingsetup}
  Euclidean map-distance readout in kilometres
\end{definition}
\begin{proof}
  This is the declared model definition of map Distance In Kilometres.
\end{proof}
\begin{definition}[map Angle At]
  \label{def:physics:phyx-mini-0690:phyxminiproblems-problemphyxmini0690-mapangleat}
  \lean{PhyXMiniProblems.ProblemPhyXMini0690.mapAngleAt}
  \uses{def:physics:phyx-mini-0690:phyxminiproblems-problemphyxmini0690-placelabel, def:physics:phyx-mini-0690:phyxminiproblems-problemphyxmini0690-lakemeetingsetup}
  The map angle first-center-second, measured in radians
\end{definition}
\begin{proof}
  This is the declared model definition of map Angle At.
\end{proof}
\begin{definition}[Matches Driving Scenario]
  \label{def:physics:phyx-mini-0690:phyxminiproblems-problemphyxmini0690-matchesdrivingscenario}
  \lean{PhyXMiniProblems.ProblemPhyXMini0690.MatchesDrivingScenario}
  \uses{def:physics:phyx-mini-0690:phyxminiproblems-problemphyxmini0690-lakemeetingsetup}
  The simultaneous-departure and simultaneous-arrival facts in the prose
\end{definition}
\begin{proof}
  This is the declared model definition of Matches Driving Scenario.
\end{proof}
\begin{definition}[Matches Supplied Figure]
  \label{def:physics:phyx-mini-0690:phyxminiproblems-problemphyxmini0690-matchessuppliedfigure}
  \lean{PhyXMiniProblems.ProblemPhyXMini0690.MatchesSuppliedFigure}
  \uses{def:physics:phyx-mini-0690:phyxminiproblems-problemphyxmini0690-angleinradiansfromdegrees, def:physics:phyx-mini-0690:phyxminiproblems-problemphyxmini0690-lakemeetingsetup, def:physics:phyx-mini-0690:phyxminiproblems-problemphyxmini0690-mapdistanceinkilometres, def:physics:phyx-mini-0690:phyxminiproblems-problemphyxmini0690-mapangleat}
  Evidence read from the supplied raster. The route for car 1 is A-L, the route for car 2 is B-L, the dashed baseline A-B is labeled 80 km, and the angle from AB to AL at A is labeled 40 degrees
\end{definition}
\begin{proof}
  This is the declared model definition of Matches Supplied Figure.
\end{proof}
\begin{definition}[Matches Problem Readouts]
  \label{def:physics:phyx-mini-0690:phyxminiproblems-problemphyxmini0690-matchesproblemreadouts}
  \lean{PhyXMiniProblems.ProblemPhyXMini0690.MatchesProblemReadouts}
  \uses{def:physics:phyx-mini-0690:phyxminiproblems-problemphyxmini0690-durationinhours, def:physics:phyx-mini-0690:phyxminiproblems-problemphyxmini0690-speedinkilometresperhour, def:physics:phyx-mini-0690:phyxminiproblems-problemphyxmini0690-lakemeetingsetup}
  Numerical physical readouts stated in the problem text
\end{definition}
\begin{proof}
  This is the declared model definition of Matches Problem Readouts.
\end{proof}
\begin{definition}[Has Physical Lake Meeting Parameters]
  \label{def:physics:phyx-mini-0690:phyxminiproblems-problemphyxmini0690-hasphysicallakemeetingparameters}
  \lean{PhyXMiniProblems.ProblemPhyXMini0690.HasPhysicalLakeMeetingParameters}
  \uses{def:physics:phyx-mini-0690:phyxminiproblems-problemphyxmini0690-durationinhours, def:physics:phyx-mini-0690:phyxminiproblems-problemphyxmini0690-speedinkilometresperhour, def:physics:phyx-mini-0690:phyxminiproblems-problemphyxmini0690-lakemeetingsetup}
  Nondegeneracy and positivity conditions for the depicted physical branch. None fixes car 2's speed or selects an answer choice
\end{definition}
\begin{proof}
  This is the declared model definition of Has Physical Lake Meeting Parameters.
\end{proof}
\begin{definition}[Satisfies Straight Route Kinematics]
  \label{def:physics:phyx-mini-0690:phyxminiproblems-problemphyxmini0690-satisfiesstraightroutekinematics}
  \lean{PhyXMiniProblems.ProblemPhyXMini0690.SatisfiesStraightRouteKinematics}
  \uses{def:physics:phyx-mini-0690:phyxminiproblems-problemphyxmini0690-durationinhours, def:physics:phyx-mini-0690:phyxminiproblems-problemphyxmini0690-speedinkilometresperhour, def:physics:phyx-mini-0690:phyxminiproblems-problemphyxmini0690-lakemeetingsetup, def:physics:phyx-mini-0690:phyxminiproblems-problemphyxmini0690-mapdistanceinkilometres}
  The governing straight-route, constant-speed kinematics for the two cars. The same duration appears in both equations because the cars leave and arrive simultaneously. No numerical value for car 2's speed occurs in these laws
\end{definition}
\begin{proof}
  This is the declared model definition of Satisfies Straight Route Kinematics.
\end{proof}
\begin{definition}[Answer Choice]
  \label{def:physics:phyx-mini-0690:phyxminiproblems-problemphyxmini0690-answerchoice}
  \lean{PhyXMiniProblems.ProblemPhyXMini0690.AnswerChoice}
  Labels of the four displayed answers
\end{definition}
\begin{proof}
  This is the declared model definition of Answer Choice.
\end{proof}
\begin{definition}[speed In Kilometres Per Hour]
  \label{def:physics:phyx-mini-0690:phyxminiproblems-problemphyxmini0690-answerchoice-speedinkilometresperhour}
  \lean{PhyXMiniProblems.ProblemPhyXMini0690.AnswerChoice.speedInKilometresPerHour}
  \uses{def:physics:phyx-mini-0690:phyxminiproblems-problemphyxmini0690-speedinkilometresperhour, def:physics:phyx-mini-0690:phyxminiproblems-problemphyxmini0690-answerchoice}
  Speed printed beside each answer choice, in kilometres per hour
\end{definition}
\begin{proof}
  This is the declared model definition of speed In Kilometres Per Hour.
\end{proof}
\begin{definition}[Is Closest Answer Choice]
  \label{def:physics:phyx-mini-0690:phyxminiproblems-problemphyxmini0690-isclosestanswerchoice}
  \lean{PhyXMiniProblems.ProblemPhyXMini0690.IsClosestAnswerChoice}
  \uses{def:physics:phyx-mini-0690:phyxminiproblems-problemphyxmini0690-speedinkilometresperhour, def:physics:phyx-mini-0690:phyxminiproblems-problemphyxmini0690-answerchoice, def:physics:phyx-mini-0690:phyxminiproblems-problemphyxmini0690-answerchoice-speedinkilometresperhour}
  A choice is closest when its displayed speed has no greater absolute error than any other displayed speed. This preserves the multiple-choice semantics without asserting that a rounded displayed value is the exact physical speed
\end{definition}
\begin{proof}
  This is the declared model definition of Is Closest Answer Choice.
\end{proof}
% --- Archon declaration topology end ---
% --- Archon physics formalization source end ---
